\documentclass[11pt]{letter}

\usepackage[margin=0.75in]{geometry}
\usepackage{setspace}
\usepackage{hyperref}
\usepackage{xcolor}
\usepackage{enumitem}
\usepackage{parskip}

\hypersetup{
    colorlinks=true,
    urlcolor=black,
    linkcolor=black
}

\setlength{\parindent}{0pt}
\setstretch{1.03}
\date{}
\begin{document}

\begin{letter}{}

\opening{Dear Hiring Committee,}

I saw the job post on the careers website and thought I'd be a good fit for it.

I am a Master's student in Computer Science at the University of Cincinnati, graduating December 2026, looking for a role as a \textbf{Software Engineer - New Grad} at \textbf{Scale AI}. With hands on experience shipping full stack web applications, integrating APIs, and building systems that process large volumes of data, I am excited about the opportunity to contribute to \textbf{Scale AI} and help move AI applications from pilot to production impact.

\textbf{Here are reasons that make me a great fit for the role:}

\begin{itemize}[leftmargin=*, itemsep=0.4em]
    \item \textbf{Product engineering and shipping customer-facing features:} I built and shipped PolicyPilot, a full stack TypeScript app (React SPA, Express REST API, and an MCP server) for high stakes eligibility decisions, and earlier a full stack React and MySQL workflow application for the Government of Maharashtra that digitized record management and improved operational efficiency by 40\%.

    \item \textbf{Prior software engineering internships and shipping at scale:} Across two internships and a full time role at Siemens Industry Software, I executed 44 production environment upgrades in 4 months, replaced legacy Ruby workflows with Python automation that cut deployment time by 40\%, and built a self service developer platform adopted by 15 engineers.

    \item \textbf{Building systems that process large volumes of data:} I built an ML system forecasting demand across 263 zones over a 2.3M row hourly dataset using a PySpark medallion pipeline, improving weighted MAPE from 18.77\% to 16.52\% and cutting RMSE by 43\%, then served predictions through a FastAPI backend with PostgreSQL and Redis caching.

    \item \textbf{Python, TypeScript, React, and hands-on LLM work:} I work daily in Python, TypeScript, React, Node.js, and SQL. My graduate research builds and evaluates agentic LLM pipelines against a published baseline (ViBR) across 53 runs on 20+ apps, achieving a 25.7x latency reduction while preserving task completion fidelity.
\end{itemize}

I am available for full-time work after December 2026, my graduation date. As a Master's student at the University of Cincinnati and someone who values professionalism and growth, I am enthusiastic about contributing to \textbf{Scale AI} and the engineering behind reliable AI systems.

Please find my CV attached below. I look forward to hearing from you. I would be excited to help build the products and infrastructure that take AI applications from pilot to real world impact.

{\setlength{\parskip}{0pt}
Regards,\\
Tanmay Bhuskute,\\
bhuskutv@mail.uc.edu $\mid$ tanmay.v.bhuskute@gmail.com $\mid$ +1 7328583027,\\
3424 Brookline Ave, Flat 19, Cincinnati, OH 45220.
}

\end{letter}

\end{document}
